2026-07-17 07:41:04.194 [info] [main] Registration level: Info

\subsection{Principios de arquitectura de software}

La arquitectura de software se define como el conjunto de decisiones estructurales significativas que determinan la organización de un sistema, sus componentes, las relaciones entre ellos y los principios que guían su diseño y evolución \cite{tecnovy_arch}. Según el enfoque contemporáneo basado en atributos de calidad, la arquitectura debe diseñarse priorizando requisitos no funcionales —como el rendimiento, la seguridad, la mantenibilidad, la disponibilidad y la escalabilidad— por encima de la funcionalidad exclusivamente \cite{arxiv_scientific_sw}, \cite{ajbas_styles}. El proceso de diseño arquitectónico es iterativo: se estiman los atributos de calidad, se evalúa el diseño frente a dichos requisitos y, en caso de no cumplirse, se procede a una fase de transformación arquitectónica hasta lograr conformidad \cite{ajbas_styles}.

\subsection{Arquitectura por capas}

La arquitectura por capas (\textit{layered architecture}) organiza el sistema en niveles jerárquicos con responsabilidades específicas, típicamente presentación, lógica de negocio y acceso a datos, favoreciendo la separación de responsabilidades y facilitando el mantenimiento del sistema \cite{lazarenko_arch}, \cite{gfg_patterns}. No obstante, este estilo puede introducir sobrecarga de rendimiento debido al paso de datos entre capas y tiende a favorecer un escalado vertical antes que horizontal, lo que limita su flexibilidad en escenarios de alta demanda \cite{vfunction_patterns}.

\subsection{Arquitecturas en pizarra, dirigida por eventos y peer-to-peer}

La arquitectura en pizarra (\textit{blackboard}) se basa en un espacio de datos compartido al cual acceden múltiples componentes especializados que contribuyen incrementalmente a la solución de un problema, siendo especialmente útil en sistemas donde no existe un algoritmo determinístico único. La arquitectura dirigida por eventos (\textit{event-driven architecture}) organiza el flujo del sistema en torno a la generación, detección y consumo asíncrono de eventos, lo que permite un alto rendimiento en el procesamiento de grandes volúmenes de información, aunque presenta desafíos relacionados con la consistencia eventual de los datos \cite{vfunction_patterns}. Por su parte, la arquitectura peer-to-peer (P2P) distribuye responsabilidades de manera simétrica entre los nodos participantes, sin depender de un servidor central, lo que favorece la resiliencia pero complejiza la coordinación y el descubrimiento de servicios.

\subsection{Arquitectura Orientada a Servicios (SOA) y Microservicios}

La Arquitectura Orientada a Servicios (SOA) descompone la aplicación en servicios débilmente acoplados que se comunican mediante protocolos estandarizados como SOAP o HTTP/XML, favoreciendo la reutilización de funcionalidad y la interoperabilidad entre sistemas heterogéneos \cite{lazarenko_arch}, \cite{opengroup_soa}. Su evolución histórica está ligada a tecnologías como XML-RPC y al surgimiento del \textit{Enterprise Service Bus} (ESB) como mecanismo para gestionar la creciente complejidad de las integraciones \cite{opengroup_soa}. La arquitectura de microservicios, considerada una evolución de SOA, propone la construcción de sistemas como conjuntos de servicios pequeños, independientemente desplegables y organizados en torno a capacidades de negocio específicas, comunicados mediante API ligeras \cite{gfg_patterns}, \cite{ijacsa_microservices}. Estudios recientes destacan que, si bien los microservicios mejoran la escalabilidad, la tolerancia a fallos y la flexibilidad operativa —tal como lo evidencian casos de empresas como Netflix, Amazon y Uber—, también introducen retos significativos en materia de consistencia de datos, gestión de API y coordinación de transacciones distribuidas, para lo cual se recomiendan prácticas como el Diseño Guiado por el Dominio (DDD) y el patrón Saga \cite{ijacsa_microservices}. En el ámbito de la salud digital, investigaciones aplicadas han comparado plataformas basadas en SOA y en microservicios para el despliegue de servicios en entornos de computación en la nube, evidenciando ventajas de los microservicios en escalabilidad selectiva frente a la mayor rigidez de los buses de servicio en SOA \cite{mdpi_ehealth}.
